
\section{The Inverse-IFEval Testbed}
\label{sec:appendix:testbed}

\subsection{Worlds}
\label{sec:appendix:worlds}

We take every IFEval item carrying at least two verifiers, giving \LabWorlds{} worlds per seed and
\LabUnits{} maintenance histories per arm across \LabSeeds{} seeds. Covers run to two instructions
(\LabStratTwoN{} histories) or three (\LabStratThreeN{}), an order of magnitude below the median
context file of \autoref{sec:growth}. We report the two strata separately wherever they disagree.
We also permutation-sample worlds under the seed, so a world index means a different item in each
seed and $(\text{seed}, \text{world})$ is the unit of analysis. A stronger model generates the
objective $o_j$ once per world, and we freeze it with the world. It cannot drift when a maintenance
protocol changes.

\subsection{Roles and the maintenance loop}
\label{sec:appendix:loop}

Three roles run the loop and only the harness measures anything. At each of \LabT{} steps the
harness samples \LabL{} task from the world, restates the objective, and draws \LabK{} of the
world's constraints to score the answer against; the executor answers under the current prompt with
comments removed; the verifiers run. A fresh maintainer instance then adds and deletes
instructions, tagging each addition with a comment as its protocol prescribes, and it carries
nothing between steps but the prompt and what its arm persists. We offer no rewording operation, so
an instruction's text is fixed when it is written and every instruction has an unambiguous age,
matching the spell definition that \autoref{sec:measurement} recovers from commit history. Maintainer
and executor are both \texttt{claude-haiku-4-5}, chosen so the maintainer is closer in capability
to the authors of \autoref{sec:field}'s files than a frontier model would be.

\subsection{Feedback regime}
\label{sec:appendix:regime}

The harness returns a complaint in plain language, naming no constraint, parameter, or taxonomy,
and it returns nothing beyond a bare verdict when a response passes. That is the censoring
\autoref{sec:theory} assumes. \autoref{tab:verifiers} lists every constraint type the world pool
instantiates beside the complaint it produces, so a reader can judge whether the maintainer's task
is well posed: the complaint names a dimension and never a parameter, so a maintainer told that a
word count is off cannot read the target off the feedback.

The arms separate only when five conditions hold jointly: vague verdicts, censored passes,
constraint identifiers withheld, a maintainer objective that does not ask for brevity, and a
horizon of at least ten steps. No arm separated until all five held. Disclose the verifier spec in
full, and the uncommented arm settles at \RegimeNullControlExcessPct{} excess size with the effect
gone entirely. We found this regime by search, and the search order licenses it: every condition
removes information that a real maintainer also does not have.

We hit two of the five as design failures rather than choosing them as parameters. Stable
constraint identifiers leak $|C|$: given ids that persist across rounds, a maintainer infers the
cover size and prunes to it with no history at all, and every arm converges. A maintainer objective
mentioning brevity makes the ratchet an artifact of the objective rather than of the information.
Dropping the minimality objective entirely left the delete-to-add rate essentially unchanged, which
rules that reading out.

\begin{table*}[t]
  \centering\small
  \begin{tabular}{@{}llr>{\raggedright\arraybackslash}p{0.44\textwidth}@{}}
    \toprule
    Family & Verifier type & Worlds & Complaint returned on failure \\
    \midrule
    keywords & \texttt{existence} & 69 & the response is missing some required content \\
     & \texttt{forbidden\_words} & 81 & the response contains wording it must not use \\
     & \texttt{frequency} & 75 & certain wording is not used the right number of times \\
     & \texttt{letter\_frequency} & 57 & a certain letter does not appear the right number of times \\
    \addlinespace
    length & \texttt{number\_paragraphs} & 48 & the response is not split into the right number of paragraphs \\
     & \texttt{number\_sentences} & 99 & the response has the wrong number of sentences \\
     & \texttt{number\_words} & 108 & the response's word count is off \\
    \addlinespace
    content & \texttt{number\_placeholders} & 42 & the response needs a different number of bracketed placeholders \\
     & \texttt{postscript} & 36 & the response is missing a postscript \\
    \addlinespace
    format & \texttt{json\_format} & 27 & the response is not valid JSON \\
     & \texttt{multiple\_sections} & 24 & the response is not organized into the expected sections \\
     & \texttt{number\_bullet\_lists} & 54 & the response does not have the right number of bullet lists \\
     & \texttt{number\_highlighted\_sections} & 87 & the response needs a different number of highlighted sections \\
     & \texttt{title} & 36 & the response is missing a properly formatted title \\
    \addlinespace
    case & \texttt{capital\_word\_frequency} & 39 & the number of fully-capitalized words is off \\
     & \texttt{english\_capital} & 42 & the response's letter casing is wrong \\
     & \texttt{english\_lowercase} & 75 & the response's letter casing is wrong \\
    \addlinespace
    punctuation & \texttt{no\_comma} & 87 & the response's punctuation breaks a rule \\
    \addlinespace
    start / end & \texttt{end\_checker} & 27 & the response does not end the required way \\
     & \texttt{quotation} & 69 & the response does not begin and end the required way \\
    \bottomrule
  \end{tabular}
  \caption{
    The hidden constraint grammar and the only signal the maintainer receives about it. Rows: the
    verifier types instantiated across the canonical run's worlds, grouped by IFEval family;
    columns: the number of $(\text{seed}, \text{world})$ pairs whose hidden constraint set contains
    that type, and the complaint the harness returns when a response fails it. The complaint names
    a dimension and never a parameter, a threshold, or the verifier's identity --- a maintainer
    told ``the response's word count is off'' cannot read the target off the feedback, which is
    the censoring \S\ref{sec:theory} assumes. Every type carries a distinct complaint except the
    two casing checks, which share one by construction. Counts sum above the world count because
    a world carries two or three constraints.
  }
  \label{tab:verifiers}
\end{table*}


\subsection{Channels}
\label{sec:appendix:channels}

An instruction is an imperative addressed to the executor, stating what the response must do and
never why. A comment is addressed to the next maintainer and carries whatever its arm persists.

The prompt has one comment syntax, and it attaches a comment to an instruction and to nothing else.
A maintained prompt is a sequence of lines, one per instruction, each written
\texttt{[d\emph{i}] \emph{text}}, where \emph{i} is an identifier the harness assigns from a
counter at addition; a \texttt{\#} and the comment follow on the same line whenever the arm carries
one. The identifier is not a position. It never changes and it is never reused, so a comment naming
\texttt{d\emph{i}} names the same instruction ten rounds later, and the executor's rendering
replaces it with a positional list. The maintainer writes a comment once, with the addition it
annotates, and cannot edit it afterwards any more than it can reword an instruction; under every
protocol whose only channel is the comment, this paper's included, the harness truncates it at
\LabCommentCap{} characters. Each addition also carries a \texttt{cause}, one phrase naming the
requirement the maintainer believes it targets, which we record and never render into any prompt.

The syntax offers nothing above the instruction line. There is no file-level comment, no preamble,
and no comment on a deletion --- the delete operation takes an identifier and carries no text at
all. A rationale not attached to a live instruction therefore has nowhere to live and dies with the
instruction it annotates, which is the constraint \autoref{eq:ratchet} formalizes rather than an
implementation shortcut: a maintainer free to annotate the prompt as a whole would keep exactly the
history this design censors, and every arm would converge.
The harness strips comments before the executor reads the prompt and rejects any instruction whose
text cites a past failure, a verification result, or a step, so no arm can route provenance through
the instruction body. \autoref{fig:prompts} gives the maintainer every word of its own instructions,
and \autoref{fig:renderings} follows one real prompt from the canonical run through both
projections. We produce the second by calling the two render functions on that run's own stored
state rather than by describing them, because the strip runs as code and not as a request the
executor could disregard.
\autoref{sec:appendix:failures} measures what the strip is worth by lifting it.

\begin{figure*}[t]
  \small
  \textbf{(a) Maintainer system prompt (all arms)}
  \begin{verbatim}
You maintain a reusable prompt (a numbered list of directives) for an assistant
  that repeatedly performs variations of one task. After each round, some
  responses were scored against hidden requirements; you see this round's scores
  below.

Optimize the prompt for LONG-RUN average score. You may add directives and delete
  existing ones.

Each response is checked against a random sample of the hidden requirements. You
  are NOT told how many requirements exist, which ones were sampled, or which
  directive (if any) satisfied a passing one; a failure may include a description
  of what went wrong.

Rules for directives:
- A directive is one imperative sentence addressed to the assistant about what the
  RESPONSE must do. It must NEVER mention scores, failures, rounds, or history.
- You cannot reword a directive; to change one, delete it and add a new one.
- With each add, supply `comment` (a note to the future maintainer: which observed
  failure prompted this, in your own words, e.g. "task 1 was flagged for comma
  use") and `cause` (one short phrase naming the requirement you believe it
  targets, e.g. "no commas allowed").

Return only the operations. An empty ops list keeps the prompt unchanged.
  \end{verbatim}
  \textbf{(b) Protocol addendum (informative-comments arm only)}
  \begin{verbatim}
Your ONLY memory across rounds is the `comment` attached to each directive (at
  most 1,024 characters, frozen when the directive is added). Write every comment
  as a commit message to all future maintainers, containing:
- "r<N>:" — this round's number, inferred as one more than the highest r<N> in the
  existing comments (r1 if there are none);
- the failure description that prompted this directive, quoted, and how many
  rounds it has been recurring;
- the lineage of this failure: which approaches were already tried against it and
  FALSIFIED (copy forward any falsifications recorded in the comments of
  directives you are deleting now — a deleted directive's comment vanishes with
  it);
- if this directive replaces a deleted one that might still have been doing work,
  the deleted directive's exact text in quotes, so a successor can restore it
  verbatim.
Before deleting any directive, read its comment: if its failure has not recurred
  since the round it was added, the directive is likely PREVENTING that failure —
  keep it. Delete a directive when its comment's failure kept recurring anyway
  (falsified), or when it duplicates another directive's target.
  \end{verbatim}
  \caption{
    What the maintainer is told, and the only thing that differs across arms. \textbf{(a)} goes to
    every arm; \textbf{(b)} goes to the informative-comments arm alone, and everything that
    protocol persists is written into the per-instruction comment at write time. No arm differs
    from another in any other input. \autoref{fig:renderings} shows what the two roles then see.
    Line breaks are set for this page.
  }
  \label{fig:prompts}
\end{figure*}

\begin{figure*}[t]
  \small
  \textbf{(a) The maintained prompt as one document (objective, instructions, comments)}
  \begin{verbatim}
Objective: Write a blog post about the sleek new magistrates.

[d1] End your response with a postscript section labeled 'P.S.'  # r1: task 0
  failed with "the response is missing a postscript"; no prior attempts
[d10] Ensure your response has between 300 and 400 words.  # r6: task 0 failed
  with "the response's word count is off"; truncated output suggests responses may
  be too short or too long. Adding explicit word count range to target this
  requirement.
[d12] Include exactly one bullet list in your response, with at least three items.
  # r7: task 0 failed with "the response does not have the right number of bullet
  lists"; this is a new failure not seen before. The response needs a structured
  bullet list to satisfy this requirement.
[d14] Include exactly three bullet lists in your response, each with at least
  three items.  # r9: task 0 failed with "the response does not have the right
  number of bullet lists"; d13 specified two lists but requirement demands three.
  Escalating from two to three lists.
  \end{verbatim}
  \textbf{(b) What the maintainer is handed each round}
  \begin{verbatim}
## Current prompt directives
(panel (a)'s commented list, verbatim)

## This round's scores
task 0: a requirement PASSED
task 0: a requirement (the response does not have the right number of bullet
  lists) FAILED
  \end{verbatim}
  \textbf{(c) What the executor is handed: the same instructions, comments stripped}
  \begin{verbatim}
Follow ALL of these directives in your response:
1. End your response with a postscript section labeled 'P.S.'
2. Ensure your response has between 300 and 400 words.
3. Include exactly one bullet list in your response, with at least three items.
4. Include exactly three bullet lists in your response, each with at least three
  items.

--- user message ---

(panel (a)'s objective, verbatim)
  \end{verbatim}
  \textbf{(d) Placebo comment pool}
  \begin{verbatim}
- added after one of the responses was flagged in an earlier round
- a response failed a requirement around this in a previous round
- several earlier rounds seemed to have trouble with this
- this came up in the scores a while back
- added to address a recurring issue seen earlier
- an earlier round's feedback pointed at something like this
- one of the tasks was marked down for this before
- kept seeing failures that looked related to this
  \end{verbatim}
  \caption{
    One real maintained prompt and the two projections of it the roles receive, from the canonical
    run's protocol arm at its last round: world 170 of seed 0, ending at 4 instructions against an
    arm median of 2, the largest prompt that fits this page. \textbf{(a)} is the artifact the
    maintainer edits and the form this paper argues for --- objective, instructions, and the
    comment each was written with, in one document that no role receives. \textbf{(b)} pairs that
    list with a censored report naming a dimension and never a parameter. \textbf{(c)} is the
    executor's call, produced by the render function rather than described: comments gone,
    identifiers gone, objective moved to the user message. The strip between them is code, not a
    request. \textbf{(d)} is the pool the placebo arm substitutes for real comments. System prompts
    are \autoref{fig:prompts}; notation is \apxref{sec:appendix:channels}'s.
  }
  \label{fig:renderings}
\end{figure*}


